\subsection{应用.}

* A) 群的酉表示：

设 E 是一个复希尔伯特空间，G 是一个群，\(\pi\) 是 G 在 E 上的一个酉表示，即从 G 到 E 的自同构群中的一个同态。设 \(E^{G}\) 是 E 中在 \(\pi(\mathbf{G})\) 下不变的所有向量所成的希尔伯特子空间。对每个 \(x \in E\)，设 \(K_{x}\) 是 x 的轨道的闭凸包。在 E 中固定一点 x。

我们将证明 \(K_{x}\) 中存在唯一的一个在 \(\pi(\mathbf{G})\) 下不变的点，即 x 在 \(E^{G}\) 上的投影。由 IV, 第 44 页 的推论（应用于 E 的底实向量空间），\(K_{x}\) 中存在一个在 \(\pi(\mathbf{G})\) 下不变的点；设 a 是这样一个点；那么 \(a \in E^{G}\)。设 P 是所有 \(y \in E\) 中满足 y - x 与 \(E^{G}\) 正交者所成的集；我们立即看出 P 闭、凸且在 \(\pi(\mathbf{G})\) 下不变；因此 \(x \in P\)，从而 \(K_{x} \subset P\)，最后 \(a \in P\)。换句话说，a - x 与 \(E^{G}\) 正交；于是 a 是 x 到 \(E^{G}\) 上的投影。

* B) 希尔伯特空间中算子的迹：

假设表示 \(\pi\) 是不可约的，也就是说，不存在 E 的异于 \(\{0\}\) 且异于 E 的希尔伯特子空间在 \(\pi(\mathbf{G})\) 下不变。设 \(\mathbf{F} = \mathcal{L}^{2}(\mathbf{E})\) 是 E 的所有 Hilbert-Schmidt 自同态所成的希尔伯特空间，赋以内积 \(\langle u|v\rangle = \operatorname{Tr}(u^{*}v)\)。我们定义从 G 到 F 中的一个酉表示 \(\lambda\)，其公式为

\[
\lambda (g). u = \pi (g) u \pi (g) ^ {- 1} \quad (u \in \mathrm{F}, g \in \mathrm{G}).\tag{3}
\]

空间 \(F^{G}\)（E 中在 \(\lambda(\mathbf{G})\) 下不变的所有元素所成的空间）由 E 的与 \(\pi(g)\) 交换（对所有 \(g \in G\)）的 Hilbert-Schmidt 自同态 u 组成。由 Schur 引理，这样的 u 是一个位似。于是我们必须考虑两种情形：

\begin{enumerate}
\def\labelenumi{\arabic{enumi})}
\item
  若 E 是无穷维的，则 \(F^{G} = \{0\}\)；
\item
  若 \(\mathrm{E}\) 是有限维的，则 \(\mathrm{F} = \mathcal{L}(\mathrm{E})\) 且 \(\mathrm{F}^{\mathrm{G}} = \mathbf{C}.1_{\mathrm{E}}\)。
\end{enumerate}

把 \(A\) ) 的结果应用于酉表示 \(\lambda\)，我们得到下面的定理：

设 \(u \in \mathcal{L}^2(\mathrm{E})\)，设 \(\mathbf{A}_u\) 是 \(\mathcal{L}^2(\mathrm{E})\) 中的一个闭凸包，即自同态 \(\pi(g) u\pi(g)^{-1}\)（\(\mathrm{E}\) 的自同态，\(g\) 取遍 \(\mathbf{G}\)）所成之集的闭凸包。若 \(\mathrm{E}\) 是无穷维的，我们有 \(0 \in \mathbf{A}_u\)。若 \(\mathrm{E}\) 是有限维的，维数为 \(d\)，则存在唯一的位似

属于 \(\mathbf{A}_u\)，即投影 \(\frac{1}{d}\mathrm{Tr}(u).1_{\mathrm{E}}\)（\(u\) 到子空间 \(\mathbf{C}.1_{\mathrm{E}}\) 上的投影，该子空间含于 \(\mathcal{L}^2 (\mathrm{E}).\)）。

\begin{enumerate}
\def\labelenumi{\Alph{enumi})}
\setcounter{enumi}{2}
\tightlist
\item
  紧群的 Haar 测度：
\end{enumerate}

设 G 是一个紧群，设 \(E = \mathcal{C}(\mathbf{G}, \mathbf{R})\) 是 G 上所有实值连续函数所成的 Banach 空间，赋以范数

\[
\| f \| = \sup _ {x \in \mathbf {G}} | f (x) |.\tag{4}
\]

对所有 \(x \in G\)，我们用下列公式定义 E 的自同构 \(\gamma_{x}\) 与 \(\delta_{x}\)

\[
\gamma_ {x} f (y) = f \left(x ^ {- 1} y\right), \quad \delta_ {x} f (y) = f (y x)\tag{5}
\]

（对 \(y \in \mathbf{G}\)、\(f \in \mathrm{E}\)）。

设 \(f \in E\)，设 \(\Gamma_f\)（相应地，\(\Delta_f\)）表示所有函数 \(\gamma_x f\)（相应地，\(\delta_x f\)）所成之集（\(x\) 取遍 G）在 E 中的闭凸包。我们将证明存在唯一的常值函数 \(\mu(f)\) 属于 \(\Gamma_f\)，唯一的常值函数 \(\mu'(f)\) 属于 \(\Delta_f\)，并且这两个常数相等。

显然，G 上的一个连续函数在 E 的自同构 \(\gamma_{x}(\text{resp. } \delta_{x})\) 下不变，当且仅当它是常值函数。现在，对 G 中的 x，所有函数 \(\gamma_{x}f(\text{resp. } \delta_{x}f)\) 所成之集在 E 中是紧的，因为映射 \(x \mapsto \gamma_{x}f(\text{resp. } x \mapsto \delta_{x}f)\)（从 G 到 E 中的映射）是连续的 (GT, X, §3, No.~4, 定理 3)。由此 (II, 第 25 页, 命题 3) 推出 \(\Gamma_{f}(\text{resp. } \Delta_{f})\) 是 E 中对于由范数定义的拓扑的紧集，从而对于 \(\sigma(\text{E}, \text{E}')\) 也是紧的。由 Ryll-Nardzewski 定理 (IV, 第 43 页, 定理 2)，\(\Gamma_{f}\) 与 \(\Delta_{f}\) 中存在常值函数。剩下的要证明：若 \(c_{1} \in \Gamma_{f}\) 与 \(c_{2} \in \Delta_{f}\) 是常数，则 \(c_{1} = c_{2}\)。

设 \(\varepsilon > 0\)。根据假设，存在点 \(x_{1}, \ldots, x_{n}, y_{1}, \ldots, y_{n}\)（属于 \(\mathbf{G}\)）与正实数 \(\lambda_{1}, \ldots, \lambda_{n}, \mu_{1}, \ldots, \mu_{m}\)，使得

\[
\lambda_ {1} + \dots + \lambda_ {n} = \mu_ {1} + \dots + \mu_ {m} = 1.\tag{6}
\]

\[
\sup _ {x \in \mathbf {G}} \left| \sum_ {i = 1} ^ {n} \lambda_ {i} f (x _ {i} x) - c _ {1} \right| \leqslant \varepsilon ,\tag{7}
\]

\[
\sup _ {x \in \mathbf {G}} \left| \sum_ {j = 1} ^ {m} \mu_ {j} f (x y _ {j}) - c _ {2} \right| \leqslant \varepsilon .\tag{8}
\]

令 \(r = \sum_{i,j} \lambda_i \mu_j f(x_i y_j)\)。那么 \(r - c_1 = \sum_{j=1}^{m} \mu_j a_j\)，其中 \(a_j = \sum_{i=1}^{n} \lambda_i f(x_i y_j) - c_1\)；由 (7)，我们有 \(|a_j| \leqslant \varepsilon\)（对 \(1 \leqslant j \leqslant m\)），从而 \(|r - c_1| \leqslant \varepsilon\)。类似地，我们证明不等式 \(|r - c_2| \leqslant \varepsilon\)，从而 \(|c_1 - c_2| \leqslant 2\varepsilon\)。由于 \(\varepsilon\) 任意，我们得到 \(c_1 = c_2\)，如所断言。

根据 \(\mu(f)\) 的定义，对每个 \(\varepsilon > 0\)，我们可以找到和为 1 的正数 \(\lambda_1, \ldots, \lambda_n\) 与元素 \(x_1, \ldots, x_n\)（属于 \(\mathbf{G}\)），使得 \(\left|\sum_{i=1}^{n} \lambda_i f(x_i x) - \mu(f)\right| \leqslant \varepsilon\) 对所有 \(x \in \mathbf{G}\) 成立。

直接可知，对 \(f, g\)（属于 \(\mathbf{E}\)）以及每个标量 \(\lambda\)，我们有 \(\Gamma_{f + g} \subset \Gamma_f + \Gamma_g\) 与 \(\Gamma_{\gamma f} = \lambda \Gamma_f\)，从而我们有关系式 \(\mu(f + g) = \mu(f) + \mu(g)\) 与 \(\mu(\lambda f) = \lambda \mu(f)\)。因此，映射 \(\mu : f \mapsto \mu(f)\)（从 \(\mathbf{E}\) 到 \(\mathbf{R}\) 中的映射）是紧空间 \(\mathbf{G}\) 上的一个均值 (IV, 第 40 页)；* 换句话说 \(\mu\) 是 \(\mathbf{G}\) 上的一个正测度，且满足 \(\mu(\mathbf{G}) = 1\)。*

直接可知 \(\mu\) 在 G 的左平移下不变，且等式 \(\mu(f) = \mu'(f)\) 蕴含 \(\mu\) 也在右平移下不变。* 换句话说，\(\mu\) 是 G 上的一个左测度且是一个右测度 (INT, VII, § 1, No.~2, def. 2)。*

\paragraph{* D) 不变测度的存在性：}

设 X 是一个 Hausdorff 拓扑空间，\(\mu\) 是 X 上的一个正有界测度，G 是 X 的同胚所成的一个群。假设对所有 \(g \in G\)，测度 \(g \cdot \mu\)（\(\mu\) 在映射 \(g : X \to X\) 下的像）以 \(\mu\) 为基。设 \(u_g\) 是 X 上的一个正的 \(\mu\)-可积函数，使得 \(g \cdot \mu = u_g \cdot \mu\)。再假设存在 X 上的两个正的 \(\mu\)-可积函数 \(\phi\) 与 \(\psi\)，它们不是 \(\mu\)-零的，且使得 \(\phi \leqslant u_g \leqslant \psi\)（\(\mu\)-几乎处处，对所有 \(g \in G\)）。我们将证明存在 X 上的一个正有界测度 \(v \neq 0\)，它以 \(\mu\) 为基且在 G 下不变。

设 P 是 Banach 空间 \(E = L^{1}(X, \mu)\) 中由满足 \(\phi \leqslant f \leqslant \psi\)（\(\mu\)-几乎处处）的函数 f 的类所成的子集。那么 P 对于弱化拓扑 \(\sigma(E, E')\) 是紧的。映射 \(h \mapsto h \cdot \mu\)（从 P 到 X 上有界实测度所成的 Banach 空间 \(F = \mathcal{M}^{b}(X)\) 中的映射）是从 P 到 E 的一个子集 \(P_{1}\) 上的一个双射，该子集对于拓扑 \(\sigma(F, F')\) 凸且紧。根据假设，\(g \cdot \mu \in P_{1}\) 对所有 \(g \in G\) 成立。设 K 是所有测度 \(g \cdot \mu\) 所成之集的闭凸包。对所有 \(g \in G\)，映射 \(v \mapsto g \cdot v\) 是 K 的一个等距仿射变换。由 Ryll-Nardzewski 定理 (IV, 第 43 页, 定理 2)，存在一个测度 \(v \in K\)，它在 G 下不变。我们有 \(\phi \cdot \mu \leqslant v\)，从而 \(v \neq 0\)。

\closechapterappendix
